\section{Modelos visión--lenguaje}
\label{sec:marco_vlm}

Un \gls{vlm} procesa imagen y texto en un mismo modelo, y hay dos formas de construirlo, que aparecieron en
ese orden. La primera, desde 2021, entrena un codificador de imagen y uno de texto para que proyecten al mismo
espacio, de modo que una imagen y su descripción queden cerca. El resultado sirve para clasificar y para
buscar casos parecidos, pero no redacta. La segunda, desde 2023, conecta un codificador visual a la entrada de
un \gls{llm}, que entonces genera texto a partir de la imagen. En ambas, el codificador de imagen suele ser un
\gls{vit} como los de la Sección~\ref{sec:marco_vit}~\cite{radford2021learning,liu2023visual}.

\subsection{Alineación contrastiva: CLIP}
\label{sec:marco_vlm_clip}

\Gls{clip} entrena un codificador de imagen y uno de texto sobre pares imagen--texto, cada uno con una
proyección final a un espacio común de embeddings. La similitud entre el embedding de una imagen ($\mathbf{u}$)
y el de un texto ($\mathbf{v}$) es el coseno del ángulo que forman,
\begin{equation}
    \cos(\mathbf{u}, \mathbf{v}) = \frac{\mathbf{u}^{\top}\mathbf{v}}{\lVert\mathbf{u}\rVert\,\lVert\mathbf{v}\rVert},
    \label{eq:marco_cos}
\end{equation}
que vale 1 si los vectores apuntan en la misma dirección, 0 si son ortogonales y $-1$ si son opuestos. En cada
lote de $N$ pares, el modelo calcula la similitud entre cada imagen y cada texto y aprende a distinguir los $N$
pares correctos entre las $N^2$ combinaciones posibles, lo que acerca los embeddings de cada par y aleja los de
pares distintos (Figura~\ref{fig:marco_clip}). El modelo original se entrenó así sobre 400 millones de pares
recogidos de internet, sin ninguna etiqueta de clase~\cite{radford2021learning}.

% Figura: entrenamiento contrastivo de CLIP. Adaptada de radford2021learning (Figura 1, panel 1).
% Se usa en sec:marco_vlm_clip, después del párrafo que define el objetivo contrastivo.

\begin{figure}[H]
    \centering
    \begin{tikzpicture}[x=1mm, y=1mm]

        % Codificador de texto (arriba) y sus embeddings T_1..T_4.
        \node[figentrada, text width=22mm, minimum height=7mm] (txt) at (-26, 44) {textos del lote};
        \node[figtarea, text width=30mm, minimum height=7mm] (enct) at (22, 44) {codificador de texto};
        \draw[figflecha] (txt.east) -- (enct.west);
        \foreach \j in {1,...,4} {
            \node[draw=figverde!70!black, fill=figverde!10, minimum width=9mm, minimum height=6mm,
                font=\footnotesize] (T\j) at (\j*10, 33) {$T_{\j}$};
        }
        \draw[figflecha] (enct.south) -- ++(0, -3.5);

        % Codificador de imagen (izquierda) y sus embeddings I_1..I_4.
        \node[figentrada, text width=18mm, minimum height=7mm] (img) at (-40, 15) {imágenes del lote};
        \node[figtarea, text width=18mm, minimum height=7mm] (enci) at (-17, 15) {codificador de imagen};
        \draw[figflecha] (img.east) -- (enci.west);
        \foreach \i in {1,...,4} {
            \node[draw=figazul!80!black, fill=figazul!12, minimum width=9mm, minimum height=6mm,
                font=\footnotesize] (I\i) at (0, 33-\i*8) {$I_{\i}$};
        }
        \draw[figflecha] (enci.east) -- (I2.west |- enci.east);

        % Matriz de similitudes; la diagonal son los pares correctos.
        \foreach \i in {1,...,4} {
            \foreach \j in {1,...,4} {
                \ifnum\i=\j
                    \node[draw=fignaranja!70!black, fill=fignaranja!30, thick, minimum width=9mm,
                        minimum height=6mm, font=\scriptsize] at (\j*10, 33-\i*8)
                        {$I_{\i}{\cdot}T_{\j}$};
                \else
                    \node[draw=figtrazo, fill=figrelleno, minimum width=9mm, minimum height=6mm,
                        font=\scriptsize, text=figtrazo] at (\j*10, 33-\i*8) {$I_{\i}{\cdot}T_{\j}$};
                \fi
            }
        }
        \node[figetiqueta, align=left, anchor=west] at (47, 17)
            {diagonal: pares\\correctos, se acercan};
        \node[figetiqueta, align=left, anchor=west] at (47, 5)
            {resto: pares\\distintos, se alejan};

    \end{tikzpicture}
    \caption{Entrenamiento contrastivo de CLIP. Adaptada de \textcite{radford2021learning}.}
    \label{fig:marco_clip}
\end{figure}


Ese espacio compartido permite clasificar sin entrenar un clasificador. Cada clase se escribe como un
texto, por ejemplo «una imagen dermatoscópica de un melanoma», y la imagen se asigna a la clase cuyo
embedding de texto le queda más cerca. Esta clasificación \emph{zero-shot}, que no usa ningún ejemplo
etiquetado de la tarea, depende de cómo se redacten esos textos: cambiar la plantilla del \emph{prompt}
cambia la exactitud~\cite{radford2021learning}.

Con etiquetas disponibles, el codificador de imagen se usa como extractor de características congelado, de dos
maneras. Un clasificador lineal entrenado sobre los embeddings, el \emph{linear probe}, mide cuánta información
sobre las clases contiene la representación~\cite{radford2021learning}. El método de los \gls{knn}, en cambio,
no entrena ningún parámetro. Con una base de $m$ casos etiquetados, para una imagen consulta $x$ busca los $k$
casos cuyo embedding $f(x_j)$ tiene mayor similitud coseno con el suyo,
\begin{equation}
    \mathcal{N}_k(x) = \operatorname*{arg\,max}_{S \subseteq \{1,\dots,m\},\ |S| = k}\ \sum_{j \in S}
    \cos\bigl(f(x),\, f(x_j)\bigr),
    \label{eq:marco_vecinos}
\end{equation}
y estima la probabilidad de cada clase $c$ como la proporción de vecinos que la tienen,
\begin{equation}
    \hat{p}_{\text{kNN}}(c \mid x) = \frac{1}{k}\sum_{j \in \mathcal{N}_k(x)} \mathbf{1}[\,y_j = c\,],
    \label{eq:marco_knn}
\end{equation}
donde $y_j$ es la etiqueta del caso $j$. La predicción es la clase de mayor proporción, el voto mayoritario,
que extiende a $k$ vecinos la regla del vecino más cercano~\cite{cover1967nearest}. Así, el espacio contrastivo sirve también de índice de casos parecidos,
que la Sección~\ref{sec:marco_rep_recuperacion} usa como evidencia para el texto.

Un modelo contrastivo solo relaciona lo que vio en sus pares de entrenamiento, y los pares de internet
casi no contienen imágenes dermatoscópicas. DermLIP, de 2025, aplica el mismo esquema a alrededor de un millón
de pares imagen--texto dermatológicos, y en clasificación \emph{zero-shot} de enfermedades de la piel supera a
los modelos biomédicos de propósito general~\cite{yan2025derm1miccv}. Un modelo contrastivo, sin embargo, no
produce texto nuevo, y para redactar hace falta conectar el codificador visual con un modelo de lenguaje
generativo.

\subsection{Modelos multimodales generativos}
\label{sec:marco_vlm_generativos}

LLaVA, de 2023, es un esquema representativo de los \glspl{vlm} generativos abiertos. Un codificador
visual del tipo de \gls{clip} extrae los rasgos de la imagen, y un proyector, en LLaVA una capa lineal, los
lleva a la dimensión de los embeddings de palabra del \gls{llm}. Los tokens visuales resultantes se concatenan
con la instrucción en la secuencia de entrada de un \gls{llm} de solo decodificador, que genera la
respuesta token a token (Figura~\ref{fig:marco_llava}). La imagen llega como una parte más de la secuencia de
entrada, sin la estructura codificador--decodificador del transformer original~\cite{liu2023visual}.

LLaVA se entrena en dos etapas. En la primera, con el codificador visual y el \gls{llm} congelados, solo se
ajusta el proyector sobre pares imagen--leyenda, para que los tokens visuales caigan en una región del
espacio que el \gls{llm} interprete. En la segunda, el \emph{visual instruction tuning}, se ajustan el
proyector y el \gls{llm} sobre conversaciones en las que se pide describir, razonar o responder preguntas
sobre una imagen~\cite{liu2023visual}. Como el \gls{llm} concentra la mayoría de los parámetros, en esa
segunda etapa se le puede aplicar \gls{lora} (Sección~\ref{sec:marco_llm_lora}) en lugar del ajuste completo,
mientras el codificador visual sigue congelado y el proyector se entrena entero.

% Figura: arquitectura de un VLM generativo tipo LLaVA. Adaptada de liu2023visual (Figura 1).
% Se usa en sec:marco_vlm_generativos, después del párrafo que presenta el esquema.
% [H]: como flotante saltaba a la página siguiente y cortaba el párrafo que venía después.

\begin{figure}[H]
    \centering
    \begin{tikzpicture}[x=1mm, y=1mm]

        % Rama visual.
        \node[figentrada, text width=18mm] (img) at (0, 0) {imagen};
        \node[figtarea, text width=22mm] (enc) at (0, 16) {codificador visual\\{\footnotesize (tipo CLIP)}};
        \node[figtarea, text width=22mm] (proj) at (0, 32) {proyector\\{\footnotesize (capa lineal)}};
        \draw[figflecha] (img) -- (enc);
        \draw[figflecha] (enc) -- (proj);

        % Rama de texto.
        \node[figentrada, text width=22mm] (instr) at (40, 0) {instrucción};
        \node[fignodo, text width=22mm, fill=figrelleno] (tok) at (40, 16) {tokenizador y\\\emph{embeddings}};
        \draw[figflecha] (instr) -- (tok);

        % Secuencia de entrada al LLM: cada grupo de tokens queda encima de la rama que lo produce, y un
        % marco agrupa sus casillas para que la flecha llegue al grupo y no a una casilla suelta.
        \foreach \k in {0,...,3} {
            \node[draw=figazul!80!black, fill=figazul!12, minimum width=4.5mm, minimum height=4.5mm]
                (v\k) at (-8.25+\k*5.5, 47) {};
            \node[draw=figtrazo, fill=figrelleno, minimum width=4.5mm, minimum height=4.5mm]
                (t\k) at (31.75+\k*5.5, 47) {};
        }
        \node[draw=figtrazo, densely dotted, rounded corners=1pt, inner sep=1.2mm, fit=(v0) (v3)] (vis) {};
        \node[draw=figtrazo, densely dotted, rounded corners=1pt, inner sep=1.2mm, fit=(t0) (t3)] (txt) {};
        \node[figetiqueta, anchor=east, align=right] at (vis.west) {tokens\\visuales};
        \node[figetiqueta, anchor=west, align=left] at (txt.east) {tokens\\de texto};
        \draw[figflecha] (proj.north) -- (vis.south);
        \draw[figflecha] (tok.north) -- (txt.south);

        % LLM y salida: los dos grupos entran como una sola secuencia.
        \node[figtarea, minimum width=66mm] (llm) at (20, 64) {LLM de solo decodificador};
        \draw[figflecha] (vis.north) -- (vis.north |- llm.south);
        \draw[figflecha] (txt.north) -- (txt.north |- llm.south);
        \node[figsalida, text width=30mm] (out) at (20, 80) {texto generado};
        \draw[figflecha] (llm) -- (out);

        \node[figetiqueta, text=fignaranja!70!black, align=left, anchor=west] at (62, 32)
            {etapa 1: se entrena\\solo el proyector};
        \node[figetiqueta, text=fignaranja!70!black, align=left, anchor=west] at (62, 64)
            {etapa 2: proyector\\y LLM};

    \end{tikzpicture}
    \caption{Arquitectura de un \gls{vlm} generativo del tipo de LLaVA. Adaptada de \textcite{liu2023visual}.}
    \label{fig:marco_llava}
\end{figure}


La misma receta se llevó después al dominio médico cambiando los datos. LLaVA-Med, de 2023, la repite con
figuras y leyendas de artículos biomédicos y se entrena en menos de quince horas con ocho
GPU~\cite{li2023llavamed}; su versión 1.5 usa Mistral como modelo de lenguaje~\cite{zeng2025mmskin}. SkinVL, de 2025, la repite con
conjuntos dermatológicos, entre ellos uno construido a partir de libros de texto de la
especialidad~\cite{zeng2025mmskin}. Phi-3.5-vision combina también un codificador de imagen con un modelo de
lenguaje de solo decodificador, en un tamaño menor~\cite{abdin2024phi3}.

Un \gls{vlm} generativo, en cualquiera de estas variantes, devuelve texto y no una distribución de probabilidad
sobre las clases. Para clasificar con él se puede buscar la clase en la respuesta generada o, sin generar,
calcular con la Ecuación~\ref{eq:marco_lm} la probabilidad que el modelo asigna al nombre de cada clase como
continuación del \emph{prompt}, normalizada por su cantidad de tokens para no favorecer los nombres cortos, y
elegir la mayor. Esos puntajes ordenan las clases, pero nada garantiza que estén
calibrados, y una respuesta generada no dice cuánto confiar en ella. Antes de pedirle un reporte hay que saber
cuándo una predicción merece afirmarse.
